%%% This file contains definitions of various useful macros and environments %%%
%%% Please add more macros here instead of cluttering other files with them. %%%

\usepackage[HTML]{xcolor}   % typesetting in color
\usepackage{cleveref}       % better references

% The hyperref package for clickable links in PDF and also for storing
% metadata to PDF (including the table of contents).
% Most settings are pre-set by the pdfx package.
\hypersetup{
    unicode,
    breaklinks,
    colorlinks,
    linkcolor=blue,
    citecolor=blue!60!white,
    urlcolor=blue,
}

% Set up formatting of bibliography (references to literature)
% Details can be adjusted in macros.tex.
%
% BEWARE: Different fields of research and different university departments
% have their own customs regarding bibliography. Consult the bibliography
% format with your supervisor.
%
% The basic format according to the ISO 690 standard with numbered references
% \usepackage[natbib,style=iso-numeric,sorting=none]{biblatex}
% ISO 690 with alphanumeric references (abbreviations of authors' names)
%\usepackage[natbib,style=iso-alphabetic]{biblatex}
% ISO 690 with references Author (year)
%\usepackage[natbib,style=iso-authoryear]{biblatex}
%
% Some fields of research prefer a simple format with numbered references
% (sorting=none tells that bibliography should be listed in citation order)
\usepackage[natbib,style=numeric,sorting=none,urldate=long]{biblatex}
% Numbered references, but [1,2,3,4,5] is compressed to [1-5]
%\usepackage[natbib,style=numeric-comp,sorting=none]{biblatex}
% A simple format with alphanumeric references:
%\usepackage[natbib,style=alphabetic]{biblatex}

% Load the file with bibliography entries
\addbibresource{bibliography.bib}

%%% Switches based on thesis type

\def\TypeBc{bc}
\def\TypeMgr{mgr}
\def\TypePhD{phd}
\def\TypeRig{rig}

\ifx\ThesisType\TypeBc
\def\ThesisTypeName{bachelor}
\def\ThesisTypeTitle{BACHELOR THESIS}
\def\ThesisTypeCSGenitive{bakalářské}
\fi

\ifx\ThesisType\TypeMgr
\def\ThesisTypeName{master}
\def\ThesisTypeTitle{MASTER THESIS}
\def\ThesisTypeCSGenitive{diplomové}
\fi

\ifx\ThesisType\TypePhD
\def\ThesisTypeName{doctoral}
\def\ThesisTypeTitle{DOCTORAL THESIS}
\def\ThesisTypeCSGenitive{disertační}
\fi

\ifx\ThesisType\TypeRig
\def\ThesisTypeName{rigorosum}
\def\ThesisTypeTitle{RIGOROSUM THESIS}
\def\ThesisTypeCSGenitive{rigorozní}
\fi

\ifx\ThesisTypeName\undefined
\PackageError{thesis}{Unknown thesis type.}{Please check the definition of ThesisType in metadata.tex.}
\fi

%%% Switches based on study program language

\def\LangCS{cs}
\def\LangEN{en}

\ifx\StudyLanguage\LangCS
\else\ifx\StudyLanguage\LangEN
\else\PackageError{thesis}{Unknown study language.}{Please check the definition of StudyLanguage in metadata.tex.}
\fi\fi

%%% Minor tweaks of style

% These macros employ a little dirty trick to convince LaTeX to typeset
% chapter headings sanely, without lots of empty space above them.
% Feel free to ignore.
\makeatletter
\def\@makechapterhead#1{
  {\parindent \z@ \raggedright \normalfont
   \Huge\bfseries \thechapter\quad #1
   \par\nobreak
   \vskip 20\p@
}}
\def\@makeschapterhead#1{
  {\parindent \z@ \raggedright \normalfont
   \Huge\bfseries #1
   \par\nobreak
   \vskip 20\p@
}}
\makeatother

% This macro defines a chapter, which is not numbered, but is included
% in the table of contents.
\def\chapwithtoc#1{
\chapter*{#1}
\addcontentsline{toc}{chapter}{#1}
}

% Draw black "slugs" whenever a line overflows, so that we can spot it easily.
% \overfullrule=1mm

%%% Macros for definitions, theorems, claims, examples, ... (requires amsthm package)

% \theoremstyle{plain}
% \newtheorem{thm}{Theorem}
% \newtheorem{lemma}[thm]{Lemma}
% \newtheorem{claim}[thm]{Claim}
% \newtheorem{defn}{Definition}

% \theoremstyle{remark}
% \newtheorem*{cor}{Corollary}
% \newtheorem*{rem}{Remark}

\newtheoremstyle{examplestyle}% name
  {6pt}% Space above
  {6pt}% Space below
  {\small}% Body font
  {}% Indent amount
  {\bfseries}% Head font
  {.}% Punctuation after head
  {.5em}% Space after head
  {}% Head spec

\theoremstyle{examplestyle}
\newtheorem{exampleinner}{Example}[chapter]

\newenvironment{example}
  {\pushQED{\qed}\begin{exampleinner}}
  {\popQED\end{exampleinner}}

\crefname{exampleinner}{example}{examples}
\Crefname{exampleinner}{Example}{Examples}

%%% Style of captions of floating objects (figures etc.)

\ifcsname DeclareCaptionStyle\endcsname
\DeclareCaptionStyle{thesis}{style=base,font=small,labelfont=bf,labelsep=quad}
\captionsetup{style=thesis}
\captionsetup[algorithm]{style=thesis,singlelinecheck=off}
\captionsetup[listing]{style=thesis,singlelinecheck=off}
\fi

%%% An environment for typesetting of program code and input/output
%%% of programs.
% This would also define the `\code` command so fuck it.
% \DefineVerbatimEnvironment{code}{Verbatim}{fontsize=\small, frame=single}

% Settings for lstlisting -- program listing with syntax highlighting
\ifcsname lstset\endcsname
\lstset{
  % language=C++,
  tabsize=2,
  showstringspaces=false,
  basicstyle=\footnotesize\tt\color{black!75},
  identifierstyle=\bfseries\color{black},
  commentstyle=\color{green!50!black},
  stringstyle=\color{red!50!black},
  keywordstyle=\color{blue!75!black}}
\fi

% Floating listings, used in the same way as the figure environment
\ifcsname DeclareNewFloatType\endcsname
\DeclareNewFloatType{listing}{}
\floatsetup[listing]{style=ruled}
\floatname{listing}{Program}
\fi

%%% The field of all real and natural numbers
\newcommand{\R}{\mathbb{R}}
\newcommand{\N}{\mathbb{N}}

%%% Useful operators for statistics and probability
\DeclareMathOperator{\pr}{\textsf{P}}
\DeclareMathOperator{\E}{\textsf{E}}
\DeclareMathOperator{\var}{\textrm{var}}
\DeclareMathOperator{\sd}{\textrm{sd}}

%%% Transposition of a vector/matrix
\newcommand{\T}[1]{#1^\top}

%%% Asymptotic "O"
\def\O{\mathcal{O}}

%%% Various math goodies
\newcommand{\goto}{\rightarrow}
\newcommand{\gotop}{\stackrel{P}{\longrightarrow}}
\newcommand{\maon}[1]{o(n^{#1})}
\newcommand{\abs}[1]{\left|{#1}\right|}
\newcommand{\dint}{\int_0^\tau\!\!\int_0^\tau}
\newcommand{\isqr}[1]{\frac{1}{\sqrt{#1}}}

%%% TODO items: remove before submitting :)
\newcommand{\todo}[1]{\noindent{\color{red!}TODO: #1}}

%%% Detailed settings of bibliography

\ifx\citet\undefined\else

% Maximum number of authors of a single work. If exceeded, "et al." is used.
%\ExecuteBibliographyOptions{maxnames=2}
% The same setting specific to citations using \citet{...}
% \ExecuteBibliographyOptions{maxcitenames=2}
% The same settings specific to the list of literature
%\ExecuteBibliographyOptions{maxbibnames=2}

% Shortening first names of authors: "E. A. Poe" instead of "Edgar Allan Poe"
%\ExecuteBibliographyOptions{giveninits}
% The same without dots ("EA Poe")
%\ExecuteBibliographyOptions{terseinits}

% If your bibliography entries are hard to break into lines, try this mode:
%\ExecuteBibliographyOptions{block=ragged}

% Possibly reverse the names of the authors with the non-ISO styles:
%\DeclareNameAlias{default}{family-given}

% Use caps-and-small-caps for family names in ISO 690 style.
\let\familynameformat=\textsc

% We want to separate multiple authors in citations by commas
% (while we use semicolons in the bibliography as per the ISO standard)
\DeclareDelimFormat[textcite]{multinamedelim}{\addcomma\space}
\DeclareDelimFormat[textcite]{finalnamedelim}{\space and~}

\fi

\newcommand{\uv}[1]{``#1''}
\newcommand{\term}[1]{\textbf{#1}}
\newcommand{\code}[1]{{\normalfont{\texttt{#1}}}}
% Object from a schema category.
\newcommand{\object}[1]{\code{#1}}
% A signature of a morphism.
\newcommand{\signature}[1]{\code{#1}}

\makeatletter
\newcounter{Paper}

% Usage: \paperChapter[<TOC title>]{<label>}{<title>}{<authors>}{<publication>}
% Provide custom TOC title if the title is too long (it can overlap with the chapter number in the TOC).
\newcommand*{\paperChapter}[5][]{%
\refstepcounter{chapter}%
% This code is straight from latex2e source.
% Well, there is just `#3` instead of `\ifx...' in the source. However, this leads to overlapping of the title and the chapter number in the TOC, so we have to fix this with a custom TOC title.
\addcontentsline{toc}{chapter}{\protect\numberline{\thechapter}\ifx&#1&#3\else#1\fi}%
\refstepcounter{Paper}%
\edef\thisdisplay{\Roman{Paper}}%
\label{\detokenize{#2}}%
\protected@write\@auxout{}{%
  \string\newlabel{#2@display}{{\thisdisplay}{\thepage}}%
}%
\begin{center}
\hspace{0pt}
\vfill

{\LARGE\bfseries Paper~{\thisdisplay}\par}

\bigbreak

{\Large\bfseries #3\par}

\bigbreak\bigbreak

#4

\bigbreak\bigbreak

#5

\vfill
\hspace{0pt}
\end{center}
}

\newcount\paperRefCount

% Prints link to the given paper. Format: "I", "II", "III", ...
\newcommand{\paperRefNumber}[1]{\hyperref[\detokenize{#1}]{\ref{\detokenize{#1}@display}}}

% Prints reference to the given paper(s). Format: "Paper I", "Paper II", "Papers I, II, III". Only the number is hyperlinked.
\newcommand{\paperRef}[1]{%
  \edef\trimmedArguments{\zap@space#1 \@empty}%
  \paperRefCount=0\relax
  \@for\paper:=\trimmedArguments\do{\advance\paperRefCount by 1}%
  \ifnum\paperRefCount=1
    Paper~%
  \else
    Papers~%
  \fi
  \paperRefCount=0\relax
  \@for\paper:=\trimmedArguments\do{%
    \advance\paperRefCount by 1
    \ifnum\paperRefCount>1
      ,\penalty0\ %
    \fi
    \hyperref[\paper]{\ref{\paper @display}}%
  }%
}

\makeatother

\newcounter{autoFootnoteLabel}
\def\lastFootnoteLabel{}

\newcommand{\publishedAs}[1]{\noindent\textbf{Published as} \fullcite{#1}}

% Usage: \labeledFootnote[<label>]{<text>}
% Creates a footnote with a specific label. If the label is not provided, it will be generated automatically (a number).
\NewDocumentCommand{\labeledFootnote}{O{} m}{%
  % Determine label
  \def\currentLabel{#1}%
  \ifx\currentLabel\empty
    \stepcounter{autoFootnoteLabel}%
    \edef\currentLabel{autofn:\arabic{autoFootnoteLabel}}%
  \fi
  % Remember for later \footref
  \global\let\lastFootnoteLabel\currentLabel
  % Create footnote
  \footnote{\label{\currentLabel}#2}%
}

% Usage: \footnoteRef[<label>]
% Creates a reference to a footnote with the given label. If the label is not provided, it will refer to the last created `\labeledFootnote`.
\NewDocumentCommand{\footnoteRef}{O{}}{%
  \def\currentLabel{#1}%
  \ifx\currentLabel\empty
    \let\currentLabel\lastFootnoteLabel
  \fi
  \hyperref[\currentLabel]{\footnotemark[\getrefnumber{\currentLabel}]}%
}

% Common stuff for all papers
\newcommand{\mffAuthor}{Department of Software Engineering, Faculty of Mathematics and Physics, Charles University, Prague, Czech Republic}

\newcommand{\footnotei}[2]{%
\mbox{%
\setbox0\hbox{#2}%
\copy0%
\hspace{-\wd0}}%
\footnote{#1}%
}

% Width to use for a figure in a fbox to make it fit the page without overflowing. Measured experimentally to a sufficient precision.
\def\framedWidth{0.999275\textwidth}

% Decides whether to include the papers. See metadata.tex for the actual switch.
% Overrides labels and refs for papers.
\def\true{true}

% Usage: \paperInput{<label>}{<path>}
% Includes a paper from the given path. All \pLabel and \pRef commands in the paper will be prefixed with the given label.
\newcommand{\paperInput}[2]{%
\ifx\includePapers\true
\newpage
\begingroup
  \def\inputPrefix{#1}%

  % This approach doesn't work for \label, because it's being redefined by someone (probably at sections or sth). Yup, macros at its finest.
  % So, we have to use the \pLabel macro below.
  % \let\oldlabel\label
  % \renewcommand{\label}[1]{\oldlabel{\inputPrefix##1}}

  % So, this doesn't work either ... gues what, if \ref is used in a caption, it's written expdanded to \oldRef in the .aux file, which is, ofc, undefined when used in LOF.
  % \let\oldRef\ref
  % \renewcommand{\ref}[1]{\oldRef{\inputPrefix##1}}

  \input{#2}
\endgroup
\fi
}

% Use this instead of \label and \ref in the included papers.
\newcommand{\pLabel}[1]{\label{\inputPrefix#1}}
\newcommand{\pRef}[1]{\ref{\inputPrefix#1}}

% Usage: \defineQuestionPrefix{<prefix>}{<marker>}
% Defines a new marker for the given question prefix. E.g., if the marker is "Q", the questions will be numbered as "Q1", "Q2", etc.
\newcommand{\defineQuestionPrefix}[2]{%
  \newcounter{question@#1}%
  \expandafter\def\csname questionprefix@#1\endcsname{#2}%
}

% Usage: \question{<prefix>}{<label>}{<human-readable label>}{<question text>}
% Creates a question
\newcommand{\question}[4]{%
  \refstepcounter{question@#1}%
  \par\medskip
  \noindent\textbf{\csname questionprefix@#1\endcsname\arabic{question@#1}}\label{#1:#2} \ \emph{#3} \ #4\par\medskip
}

% Usage: \questionRef{<prefix>}{<label>}
% References a question
\newcommand{\questionRef}[2]{%
  \hyperref[#1:#2]{%
    \csname questionprefix@#1\endcsname\ref{#1:#2}%
  }%
}

\newcounter{publication}
\newcommand{\publication}[1]{%
\refstepcounter{publication}%
\noindent\makebox[2em][l]{\textbf{\arabic{publication}.}}%
\hangindent=2em\hangafter=1\relax%
\fullcite{#1}\par
\vspace{1em}
}

% Settings for the alltt environment (like, they are global, but we use them only there).

\newcommand{\keyword}[1]{\textbf{#1}}

\definecolor{sql}{HTML}{9673A6}
\newcommand{\sql}[1]{\textcolor{sql}{\bfseries #1}}

\definecolor{cypher}{HTML}{6C8EBF}
\newcommand{\cypher}[1]{\textcolor{cypher}{\bfseries #1}}

\definecolor{xquery}{HTML}{82B366}
\newcommand{\xquery}[1]{\textcolor{xquery}{\bfseries #1}}

\newcounter{contribution}[section]
\newcommand{\contribution}[1]{%
  \refstepcounter{contribution}%
  \paragraph{\arabic{contribution}}\emph{#1} \ %
}
  % \par\medskip
  % \noindent\textbf{\arabic{contribution}} \ \emph{#1} \ #2\par\medskip
